\section{Closing remarks and research agenda}
\label{sec:closing}

This paper has walked through the full pipeline implemented by
\code{jump-diffusion-estimation}: an exactly-solvable SDE whose only
approximation is the Bernoulli treatment of the jump count
(Section~\ref{sec:model}); five jump families behind a density-only
extension contract, with closed-form Gaussian convolutions exploited
wherever the mathematics allows (Section~\ref{sec:distributions}) and an
FFT scheme, ported from \citet{Ospina2009}, wherever it does not
(Section~\ref{sec:fft}); maximum likelihood driven by either a local
quasi-Newton method or the differential evolution configuration whose
robustness was the thesis' applied finding (Section~\ref{sec:estimation});
three inference routes of increasing robustness and cost
(Section~\ref{sec:inference}); a parametric-bootstrap likelihood-ratio
test that sidesteps the boundary and Davies pathologies by simulating the
null distribution outright (Section~\ref{sec:jumptest}); and
estimation-error-aware goodness-of-fit comparison across jump families
(Section~\ref{sec:comparison}).

A recurring pattern deserves to be named, because it \emph{is} the
library's design philosophy:

\begin{itemize}
  \item \textbf{Closed forms first, generic numerics as fallback.} Every
        numerical fallback (FFT convolution, inverse-CDF sampling) is
        validated against the closed forms available in special cases.
  \item \textbf{Pathologies are mapped to explicit poison values} ---
        $+\infty$ objectives, \code{NaN} standard errors --- never
        silently absorbed.
  \item \textbf{When asymptotics are doubtful, simulate.} The bootstrap
        appears three times (intervals, jump test, KS p-value), each time
        replacing a questionable asymptotic calibration with an honest
        Monte Carlo one.
  \item \textbf{Reproducibility is an argument, not an accident:} every
        stochastic routine accepts a seed and derives per-replicate seeds
        deterministically from it.
\end{itemize}

\subsection*{Open questions}

The library makes a set of methods available; it does not yet certify
their finite-sample behavior. The companion research program, for which
this library is the experimental engine, targets exactly that gap:

\begin{enumerate}
  \item \textbf{Coverage.} Empirical coverage of the Wald, profile and
        bootstrap intervals across jump families, sample sizes, jump
        intensities and jump-to-diffusion scale ratios --- when does
        profiling recover the nominal coverage that Wald loses?
  \item \textbf{Jump test calibration.} Size and power of the bootstrap
        LR test of Section~\ref{sec:jumptest} against the naive
        $\chi^2$ and the boundary-corrected mixture calibrations.
  \item \textbf{Identifiability.} Mapping the region of
        $(\lambda\Delta t,\ \text{jump scale}/\sigma\sqrt{\Delta t})$
        where $p$, $\sigma$ and the jump scale can be jointly recovered,
        via the conditioning of the observed information and the sampling
        correlation of $(\hat p, \hat\sigma)$.
  \item \textbf{Bernoulli error.} Quantifying the bias of
        Decision~\ref{dec:bernoulli} against the exact Poisson mixture as
        $\lambda\Delta t$ grows.
  \item \textbf{Optimizer geography.} Where in parameter space multistart
        L-BFGS-B fails and differential evolution succeeds, extending
        \citet{Ardia2011} from a demonstration to a map.
\end{enumerate}

Each question already has its instrument in the library; what remains is
the systematic experiment. That study is deliberately kept out of this
document: here we fixed \emph{what the code computes and why}, so that the
experiments --- and any reader's own use of the library --- stand on
explicit mathematical ground.
